\documentclass[A4,10pt,fleqn,dvipdfmx]{jarticle}
% fleqnオプション・・・数式の左寄せ
% dvipdfmxオプションを指定しないとpng, jpg等が読み込めない

\renewcommand{\normalsize}{\fontsize{10.4pt}{13.8pt}\selectfont}     %% 10pt x 1.04+α TeXの日本語フォントは少しだけ大きいため
\newcommand\thefontsize[1]{{#1 The current font size is: \f@size pt\par}}

\usepackage[left=20mm,right=20mm,top=25.5mm,bottom=28.5mm,footskip=10mm,headsep=19pt]{geometry}


\usepackage{layout} % 余白レイアウトを表示
\usepackage{fancyhdr} % ヘッダの設定
\pagestyle{fancy}
\fancyhead[R]{\thepage ~~/~~\pageref{lastpage}}
\fancyhead[L]{}
\fancyfoot[L]{\begin{tabular}{l}\hline {\fontsize{8.48pt}{12pt}\selectfont 2024 年度日本冷凍空調学会年次大会講演論文集（2024.9.4-6，福岡）}\end{tabular}}
\fancyfoot[C]{}
\fancyfoot[R]{}
\renewcommand{\headrulewidth}{0pt} % Turn off the bar

\usepackage{amsmath} % fleqnオプションを使用するために必要
\usepackage{txfonts}
\usepackage[deluxe]{otf} %明朝太字を使えるようにする
\usepackage[dvipdfmx]{graphicx}%この設定をすることでPS経由でなくdvipdfmxで直接PDF出力できる！！
\usepackage{caption}  \captionsetup{labelsep=period}
\usepackage{titlesec}
\titleformat{\section}[hang]{\bfseries}{\thesection. }{0pt}{}
\titleformat{\subsection}[hang]{}{\thesubsection. }{5pt}{}
\titleformat{\subsubsection}[hang]{}{\thesubsubsection}{5pt}{}
\titlespacing{\section}{0pt}{14pt}{0pt}
\titlespacing{\subsection}{0pt}{14pt}{0pt}
\titlespacing{\subsubsection}{0pt}{14pt}{0pt}	

% フロート（図・表）の間隔の設定（ 適宜調整してください ）
\setlength{\textfloatsep}{6pt}
\setlength{\intextsep}{20pt}
\setlength{\floatsep}{-40pt}

% 任意のパッケージ読み込み
\usepackage{tabularx} % 任意幅の表
\usepackage{xspace} % 英単語の後にはスペースを自動で追加
\usepackage{svg} %svg形式のグラフィックを貼るのに必要
% \usepackage{rotating}

% 予稿用の設定
\renewcommand{\figurename}{Fig.~}
\renewcommand{\tablename}{Table~}
\renewcommand{\refname}{REFERENCES}

% 参考文献リストにBibTeXを使用するための設定

% リストの項目間のスペースを-3ptに調整
\let\oldbibitem\bibitem
\renewcommand{\bibitem}{\setlength{\itemsep}{-3pt}\oldbibitem}


% natbibパッケージを使用する場合は53\UTF{FF5E}73行のコメントを外す
% \usepackage[numbers]{natbib}
% % \setlength{\bibsep}{0.5zh}
% \setlength{\bibsep}{0zh}
% \setlength{\bibindent}{0pt}  % natbib は番号方式の場合\bibhang を使わない．\bibindent を設定し，さらに下のようにnabib.sty中の定義を変更する．\itemindentも変更する必要があるが，それは\thebibliography 環境内で行う．jecon.bstを使う場合はその中で定義する．

% % 以下，natbib.sty の一部を再定義
% \makeatletter
% \renewcommand\NAT@bibsetnum[1]{\settowidth\labelwidth{\@biblabel{#1}}%
%   %  \setlength{\leftmargin}{\bibindent}\addtolength{\leftmargin}{\dimexpr\labelwidth+\labelsep\relax}%
%    \setlength{\leftmargin}{18pt}
%    \setlength{\itemindent}{-\bibindent}%
%    \setlength{\listparindent}{\itemindent}
% \setlength{\itemsep}{\bibsep}\setlength{\parsep}{\z@}%
%    \ifNAT@openbib
%      \addtolength{\leftmargin}{\bibindent}%
%      \setlength{\itemindent}{-\bibindent}%
%      \setlength{\listparindent}{\itemindent}%
%      \setlength{\parsep}{0pt}%
%    \fi
% }
% \makeatother


% 自分用のマクロ等
\newcommand{\degC}{$^\circ$C\xspace}
\renewcommand{\eqref}[1]{Eq.~(\ref{#1})\xspace}
\newcommand{\tblref}[1]{Table~\ref{#1}}
\newcommand{\figref}[1]{Fig.~\ref{#1}}

%%%%%% 原稿記述部分 スタート %%%%%%%%%%%%%%%%%%%%%%%%%%%%%%%%%%%%%%%%%%%%%%%%%%%%%%%%%%%%%%%%%%%%%%%%%%%%%%%%%%%%%%%%%%%%
\begin{document}
% \layout   % 余白設定を出力したい場合はコメントを外す
\setlength{\mathindent}{28pt}

\begin{center}
\fontsize{12.72pt}{16pt}\selectfont 

\textbf{和文表題講演論文の書き方}

\textbf{第2報：必要に応じて副題}


\vspace{14pt}

\normalsize
\textbf{Instructions for Manuscript Preparation}

\textbf{2nd Report: Subtitle as Necessary}


\vspace{15pt}

○冷凍 太郎*（日本橋大工），空調 花子**（中央工大理），食品 二郎***（大富大工）

○Taro REITO*, Hanako KUCHO** and Jiro SHOKUHIN***

*Department of Mechanical Engineering, Nihonbashi University, Chuo-ku, Tokyo, 103-0011, Japan

**Department of Physics, Chuo Institute of Technology, Meguro-ku, Tokyo, 152-0033, Japan

***Department of Chemical Engineering, Otomi University, Shibuya-ku, Tokyo, 151-0053, Japan


\end{center}

\vspace{16pt}
Abstract should be written with 10pt Times New Roman and single-space from here. The length of the abstract is limited within 150 words. 

PDF(.pdf)ファイルに変換してから提出する．PDF変換の設定は高解像度High quality を推奨するが，ファイルサイズは5 MB以下とする．原稿内の和文フォントは明朝体を，英数字は半角文字としTimes New Roman を推奨する．提出期限：2024年7月5日（金）

\vspace{4pt}
\begin{center}
Keywords: Keyword1, Keyword2, Keyword3, Keyword4, Keyword5
\end{center}

\vspace{3pt}



\section{\textbf{原稿作成要領}}
原稿はA4用紙2ページ\UTF{FF5E}6ページとし，本文は10pt明朝体を用いて1段組みで作成する．1ページあたり全角48文字×51行とする．

\subsection{原稿用紙}
\subsubsection*{(1)  用紙サイズ}
A4用紙とする．

\subsection*{(2)  余白}
上下左右に以下の寸法で余白を設定する．上部余白にはヘッダーを設け，「ページ番号／全ページ数」を記入する． 1ページ目の下部余白にはフッターを設け，「2024 年度日本冷凍空調学会年次大会講演論文集（2024.9.4-6，福岡）」の文字を8pt明朝体で記入する．

上余白：25 mm　下余白：25 mm

左余白：20 mm　右余白：20 mm

上部ヘッダー：用紙端から15 mm

下部フッター：用紙端から9 mm

\section{原稿書式}
Wordのスタイル機能を使用すると容易に書式を設定できる．

\subsection{表題}
和文表題は12pt明朝体強調文字，中央揃えとする．必要に応じて，副題を同じ書体で改行して記す．

英文表題は，書体10pt Bold Times New Romanとし，前置詞，冠詞を除く単語の先頭文字のみを大文字とする．副題は同じフォントを使用し，副題の先頭文字のみを大文字にして記す．

\subsection{著者}
副題から1行あけて，著者を記入する．書体は，和文では10pt 明朝体，英文では10pt Times New Romanとする．和文，英文ともに，講演者が35歳以上の場合は○を，35歳未満の場合は◎を講演者名の前に記入する．（35歳未満の場合，優秀講演賞選考対象者となる．）英文による著者名はFirst name FAMILY NAMEの順に記し，姓（FAMILY NAME）は大文字のみを用いる．その下に英文による所属を加える．全ての著者の所属，住所を例にならい記入する．

\subsection{Abstract講演概要}
著者所属から2行あけて，Abstractを150語以内で記入する．書体は10pt Times New Roman，両端揃えとする．

\subsection{Keywordsキーワード}
Abstractから1行あけて，英文keywordを3から5つ，中央揃えで記入する．Keyword の先頭文字のみを大文字とする．

\subsection{本文}
Keywordsから2行あける．本文は1段組とする．設定の詳細は「1．原稿作成要領」に後述する．
%後述？

\subsection{見出し}
章，項，節の見出しは以下の通りとする．見出しの前は1行あける．見出しが続く場合はその限りでない．書体は本文と同じ10pt明朝体であるが，章題のみ強調文字とする．

\subsection{式}
式と記号は，10pt英文フォント斜体字で記述する．ただし，添字は標準体とする．式の前後には1行あけ，左側2文字の字下げとする．式番号を右端に記す．本文中では，Eq. (1)とする．

\begin{eqnarray}
  G_{\rm in} -  F_{\rm in} = p_{\rm in}V_{\rm in} 
\end{eqnarray}

\subsection{図，表，写真など}
図などを配置する場合，上下左右の余白にはみ出さないようにする．図などの説明は英文とし，本文中では，Fig.1，Table 1などと記す．図，表の題目の上下を1行あける．Fig. 1，Table 1は図，表の例である．

図のみカラーの使用を認めるが，白黒印刷でも内容の把握が可能であるよう配慮する．

\begin{figure}[htbp]
  \centering
    \includegraphics[width=98.27mm]{fig1.png}
    \caption{Deviation of measured freezing point temperatures from correlations by Tillner-Roth\cite{Tillner-Roth1997}.}
\end{figure}

\clearpage % Wordテンプレートと同じ見た目になるように（本来は不要）

\begin{table}[htbp]
  \begin{center}
  \caption{Historical data of vapor pressure of ammonia.}
  \begin{tabular}{llllr}
    \hline
    Author	& Year	& $T$/K	& $p$/MPa	& Ref. \\
    \hline
    Keyes	& 1918	& 240-398	& 0.103-9.96	& 1 \\
    Cragoe	& 1920	& 195-343	& 0.00563-3.31	& 25 \\
    McKelvey	& 1923	& 195-195	& 0.1	& 26 \\
    Beattie	& 1930	& 303-405	& 1.17-11.3	& 8 \\
    Overstreet	& 1937	& 176-242	& 0.0008-0.1114	& 27 \\
    \hline
  \end{tabular}
  \end{center}
\end{table}

\vspace{-30pt} % Wordテンプレートと同じ見た目になるように（本来は不要）
\subsection{参考文献}
文献は，本文中で[1], [2,3], [4-7]のように番号で示す．

［例］Keyesら[1] はアンモニア(NH3)の飽和蒸気圧…

参考文献の一覧は英文表記にて文末に記載する．記載項目は，本文文末の例を参照のこと．

\section{提出方法}
PDFファイルに変換し，年次大会ホームページ（https://www.jsrae-nenji.org/nenji2024/）を通じて提出する． PDF変換の設定は高解像度High quality を推奨するが，ファイルサイズは5MB以下とする．ファイル容量が上限を超えた場合は，実行委員会でファイル容量を圧縮する場合があることに留意する．なお，ファイル名は，自動的に［講演申込番号］.pdfとなります．

\section{原稿提出期限}
2024年7月5日（金）必着（締切厳守）

\section{問合せ先}
\subsection{日本冷凍空調学会2024年度年次大会実行委員会}
E-mail：jsrae24@jsrae.or.jp
\subsection{日本冷凍空調学会事務局}
〒103-0011 東京都中央区日本橋大伝馬町13-7 日本橋大富ビル5階

TEL/FAX：03-5623-3223/3229

\section*{謝辞}
（科研費番号など）…　　　　　　　　　（本文終了）

\section*{NOMENCLATURE}
\begin{tabular}{lcl}
$c_p$ & : & specific heat at constant pressure, kJ$\cdot$kg$^{-1}$$\cdot$K$^{-1}$ \\
$h$ &	: & specific enthalpy, kJ$\cdot$kg$^{-1}$ \\
$p$ &	: & pressure, kPa \\
$T$ &	: & temperature, K \\
$t$ &	: & Celsius temperature, ℃ \\
\end{tabular}


% BibTeXを使わない場合は以下のコメントアウトを外す
% \begin{thebibliography}{xxx}

%   \harvarditem[{Keyes and R.B. Brownlee}]{{WKeyes and R.B. Brownlee}}{1918}{Keys2018}
%   {F.G. Keyes and R.B. Brownlee}, {The Vapor Pressure of Liquid Ammonia up to the Critical Temperature. [Part II.]}, 
%   Journal of the American Chemical Society, 40(1) (1918), pp. 25\UTF{2013}53.
  
%   \harvarditem{Tillner-Roth et al.}{1997}{Tillner-Roth1997}
%   {R. Tillner-Roth, J. Li, A. Yokozeki, H. Sato and K. Watanabe}, {Thermodynamic Properties of Pure and Blended Hydrofluorocarbon (HFC) Refrigerants}, JSRAE, (1997).

%   \harvarditem{Miyazaki and Oguchi}{2001}{Miyazaki2001}
%   {T. Miyazaki and K. Oguchi}, Proc. 2001 JSRAE Annual Conference, JSRAE, (2001), pp. 181\UTF{2013}184. (in Japanese)

%   \harvarditem{JSRAE}{2010}{JSRAE2010}
%   Technical report of JSRAE, JSRAE Corp., (2010. 4).

%   \harvarditem{EWEA}{2015}{EWEA2015}
%   Wind in power: 2015 European statistics, available from <http://www.ewea.org/statistics/>, (accessed on Apr. 14, 2016).

% \end{thebibliography}



% BibTeXを使って参考文献リストを出力
\nocite{*}  % .bibファイル中の全ての文献を参照する
\bibliographystyle{JSRAE}
\bibliography{references}

\normalfont
\section*{LaTeXテンプレート利用者に向けた注釈}
\paragraph*{概要}本テンプレートは2024年度日本冷凍空調学会年次大会講演論文集に対応したLaTeXテンプレートです．MS Word版テンプレートと可能な限り同じ体裁となるよう調整されています．プリアンブル中の書式設定および本文中の余白の設定（\verb|\vspace{..}|）は変更しないでください．
\paragraph*{参考文献リスト} \cite{Higashi2016}以降の例を新たに追加しています．具体的には，MS Word版では例示されていない，ページ番号がなく論文番号（Paper ID）が付されたプロシーディングス\cite{Higashi2016}，著者数が非常に多い場合\cite{IPCC2013}, webサイト引用で著者名がある場合\cite{ASHRAE2019}，博士論文\cite{Reito2017}および修士論文\cite{Kucho2008}の例を参考までに追加しています．



\label{lastpage} %全体のページ数をカウントする．\end{document}直前に置く．

\end{document}
